\section{The Crystallographic and Spinor Requirement}

This chapter is the second half of the selection argument. The dimension ladder narrowed the substrate to the
four-dimensional family. Within that family, one criterion remains, and it is decisive. Only $\{3,4,3,4\}$
carries a directional coin that is both crystallographic and able to host a spinor. Everything else in the model
is a scalar, so the coin is the only place where spin can hide, and the coin of $\{3,4,3,4\}$ is the only one
across all dimensions that holds it.

\subsection{Crystallographic means gauge is possible}

A coin carries a root system, and therefore a gauge group, only if it is \emph{crystallographic}. The
crystallographic restriction permits the angles built from $3$, $4$, and $6$, and forbids five-fold and
seven-fold symmetry. This is the same restriction that governs ordinary crystals. The compact regular honeycombs
in every dimension all contain a $5$ in their symbol, placing them in the icosahedral family. They are
non-crystallographic, so they have no root system, no gauge group, and no spinor. Compactness and a gauge group
are at odds among the compact regulars, and the substrate cannot have both.

\subsection{Only the 24-cell carries the spinor}

Among the crystallographic coins, the $24$-cell stands alone. Its symmetry is the $D_4$ root system, and its
direction set decomposes as $8v + 8s + 8c$, the vector and the two spinors. A cubic coin gives the vector only,
with no spinor at all. The $24$-cell appears as the cell of $\{3,4,3,4\}$, so the substrate hands over the $D_4$
root system, and with it the two spinor representations, as part of its own geometry.

The reason a smaller coin cannot carry a spinor is a theorem rather than an accident, and the cleanest case to
state it on is $\{5,3,4\}$, whose coin has $12$ directions.

\begin{theorem}[Permutation characters carry no spinor]
The $12$ face directions of the $\{5,3,4\}$ coin form a permutation representation of the icosahedral rotation
group. A permutation representation of a rotation group contains no spinor, because spinors live only in the
double cover. The relevant double cover here is the binary icosahedral group of order $120$, not the order-$60$
rotation group. The $\{5,3,4\}$ coin is therefore provably spin-$0$.
\end{theorem}

The force of the theorem is that it leaves no room to negotiate. A permutation representation simply cannot
contain a spinor, so no clever reading of the $12$ directions will produce one. The same logic applies to the
rest of the model by default. The tone and any memory are single ternary values, each one a scalar that returns
to itself under a full $2\pi$ rotation. So the substrate is Klein-Gordon-like, that is, spin-$0$, unless spin is
hiding in the directional coin. That is precisely why the coin is the only place to look, and why the $D_4$
decomposition $8v + 8s + 8c$ with its two spinors is the deciding fact.

\subsection{Spin is the sole discriminator}

To test whether spin is genuinely the discriminator, we run the identical battery on $\{5,3,4\}$ as a control.
Almost everything ports. The rule ports. Conservation ports. Gliders and motion port. The holographic $1/r^2$
correlator ports. Cosmology ports. The self-nesting tower ports. The two honeycombs agree across the whole sweep
with one exception.

\begin{result}[The $\{5,3,4\}$ control]
Run on $\{5,3,4\}$, every framework result matches $\{3,4,3,4\}$ except spin. The $12$ directions of $\{5,3,4\}$
are icosahedral, decomposing as $1 + 3 + 3' + 5$ with no spinor, so the substrate cannot host fundamental
fermions. It is not a root system, so there is no $SO(8)$ or $SO(10)$, hence no weak mixing angle
$\sin^2\theta_W = 3/8$ and no three generations. Its two-dimensional cusp gives only logarithmic gravity and
anyonic statistics. The single axis on which $\{3,4,3,4\}$ wins is spin. \cite{pollard2026vibetest}
\end{result}

This is what makes the control decisive. Spin is not one advantage among many. It is the one place the two
honeycombs differ, which means the spinor coin is the lever that forces the four-dimensional $\{3,4,3,4\}$ and
nothing else does.

\subsection{The cost of imposing spin by hand}

One could try to put spin onto $\{5,3,4\}$ by hand, and the cost of doing so is the clearest argument against it.
Three additions would be required. First, the binary icosahedral double cover, with its central element of order
two acting as $-1$. Second, a two-component spinor on every cell, which is new state the base did not have.
Third, a coin operator together with a map from the faces to the gamma matrices. That face map is non-unique.
Different choices give inequivalent theories, and there is no principle in the model to pick among them. The
trouble runs deeper still. A spinor is by nature a continuous complex $SU(2)$ object, so imposing one either
breaks the discrete-only commitment of the base or demands an awkward discretisation that has to keep the
double-cover $-1$ exact. By contrast $\{3,4,3,4\}$ hands over the spinor for free, inside the discrete $D_4$ coin,
with nothing added. One caveat belongs here. There is a topological route to spin, where a self is a
hopfion and the hopfion behaves as a fermion, and that route works on either honeycomb without the coin. So the
precise claim is narrow and exact. The coin spinor requires $\{3,4,3,4\}$. It is not claimed that $\{3,4,3,4\}$
is required for any spin whatsoever.

\subsection{The unique solution}

Collecting the requirements gives a clean intersection. The substrate must be four-dimensional, for a
three-dimensional physical space. It must be crystallographic, for a gauge group. It must carry the $D_4$ coin,
for spinors. And it must be hyperbolic, for gravity, holography, and expansion.

\begin{theorem}[The selection]
The substrate satisfying all four requirements, four-dimensional with a three-dimensional cusp, crystallographic,
$D_4$-coined, and hyperbolic, is uniquely $\{3,4,3,4\}$. It is paracompact, so its cusp is the three-dimensional
physical space. The genuinely compact four-dimensional regulars all contain a $5$, so they are
non-crystallographic and fail both the gauge and the spinor requirements. The flat twin $\{3,4,3,3\}$ has the
coin but lacks curvature, so it fails the gravity requirement.
\end{theorem}

A heuristic rubric makes the same point with a score. Weighting the spinor at $3$, crystallographic at $2$,
dimension closeness at $3$, curvature at $1$, and adding a compactness term, every tessellation can be ranked.
$\{3,4,3,4\}$ tops the list at $9.5$. The flat $D_4$ honeycomb $\{3,4,3,3\}$ is second at $9.0$, missing only the
curvature point. Every other tessellation loses either the spinor coin or the dimension. $\{3,4,3,4\}$ is the only
one that scores non-zero on every factor at once.


The phrasing changes from a choice to a derivation. The sentence "we chose $\{3,4,3,4\}$" becomes "the
requirements are met by only $\{3,4,3,4\}$ across all dimensions". The controls do double duty. They confirm that
the framework is generic, since it ports to every substrate, and they confirm that the physics is specific to
$\{3,4,3,4\}$, since the spinor coin appears nowhere else.
